% The six registration documents: the study each covers, when its data were collected, and what
% it registers. Design facts from the registrations; numbers only through generated/numbers.tex.
% See supplement/tables/README_suppA.md.
\begin{table}[tbp]
\centering
\small
\caption{\textbf{The six registration documents.} All six were filed together on the Open Science Framework (OSF) on 8 October 2026, after every study's data had been collected. Studies are listed in the paper's order; the order in which they were run was 2, 4, 1, 3, 5. \emph{Confirmatory}: the small set of tests that a registration names as the study's confirmatory tests, judged together after Holm's correction. Bounds are in rating points on 1-to-5 scales.}
\label{tab:suppA-registrations}
\begin{tabular}{@{}>{\raggedright\arraybackslash}p{3.0cm}>{\raggedright\arraybackslash}p{3.1cm}>{\raggedright\arraybackslash}p{\dimexpr\textwidth-6.1cm-4\tabcolsep\relax}@{}}
\toprule
Document & Study, data collection and analyzed sample & What it registers\\
\midrule
Study~1, short registration & Study~1; 25 February to 3 March 2025; $N$~=~\SOneN & H1, the AI label lowers both creativity and effort; H2, the drop is larger on effort; H3, the drop is the same for images with and without a human figure; H4, under the AI label ratings vary more between people and less within a person. Four confirmatory tests (two for H1, one each for H2 and H4); H3 decided by its own rule against $\pm$\SOneHThreeBound. Five word-meaning tests on the explanations as an exploratory set, and the two text checks.\\
\addlinespace
Study~1, long registration & Study~1 & The full plan, which governs where the short registration is silent: why the predictions matter for the law (section~2), design, exclusions, every model and decision rule, the robustness checks, the word lists and text analyses (section~9), which decisions were taken with the estimates already known (section~8), and the materials (section~10).\\
\addlinespace
Study~2 & Study~2; 23 November to 5 December 2024; $N$~=~\STwoN & H1, ratings of the second work rise from None to Slight to Dramatic; H2, the label's effect, averaged over levels, lies inside $\pm$\SourceBound{} on each measure; H3, whether the label's effect shrinks from None to Dramatic (no direction predicted). Twelve confirmatory tests, four per measure; ordinal checks; the two text checks on creativity explanations.\\
\addlinespace
Study~3 & Study~3; 4 to 9 March 2025; $N$~=~\SThreeN & The label, the stated time and their interaction on creativity and on effort (six confirmatory $F$~tests); an equivalence test of the Human-minus-AI gap at 100 hours against $\pm$\SourceBound; a rule deciding whether time closes the gap, the gap persists, or the result is inconclusive; ordinal checks; the two text checks on both ratings.\\
\addlinespace
Study~4 & Study~4; 27 January to 2 February 2025; $N$~=~\SFourN & H1 and H2 as in Study~2; H3, whether the hours stated for the original change ratings of the second work, alone or with the level (no direction predicted). Fifteen confirmatory tests, five per measure. Secondary: the label's effect at each level, its change with stated hours, an equivalence test of 100 hours against under 1 hour, and a refit without Human-label participants who wrote that a machine made the work; the two text checks on creativity explanations.\\
\addlinespace
Study~5 & Study~5; 21 October to 2 November 2025; $N$~=~\SFiveN & The effect of stated editing hours, and its interaction with the label, on each measure (six confirmatory $F$~tests); an equivalence test of editing time against $\pm$\SFiveEditTimeBound; the label's effect at each level against $\pm$\SourceBound, nine tests each judged at .05 without correction; ordinal checks; a second analysis with every item rescaled from 0 to 1; the two text checks on creativity explanations.\\
\bottomrule
\end{tabular}
\end{table}
